% Generated draft: requires mathematical and source review.
\documentclass[11pt]{article}
\usepackage[margin=25mm]{geometry}
\usepackage{fontspec}
\setmainfont{Latin Modern Roman}
\usepackage{amsmath,amssymb,mathtools,unicode-math}
\setmathfont{Latin Modern Math}
\usepackage{enumitem,needspace,xurl,hyperref}
\setlength{\parindent}{0pt}
\setlength{\parskip}{5pt}
\setlength{\emergencystretch}{3em}
\newcommand{\N}{\mathbb{N}}
\newcommand{\Z}{\mathbb{Z}}
\newcommand{\Q}{\mathbb{Q}}
\newcommand{\R}{\mathbb{R}}
\newcommand{\C}{\mathbb{C}}
\newcommand{\F}{\mathbb{F}}
\newcommand{\A}{\mathbb{A}}
\newcommand{\PP}{\mathbb{P}}
\newcommand{\E}{\mathbb{E}}
\newcommand{\eps}{\varepsilon}
\newcommand{\dd}{\,\mathrm{d}}
\begin{document}
\section*{UnsolvedMath ID: 30003336}
\textbf{Source problem code:} catalogue entry
\subsection{Definitions and mathematical statement}
For a graph \(H\), write \(V(H)\) and \(E(H)\) for its vertex and edge sets, and let \(\delta(H)=\min_{v\in V(H)}\deg_H(v)\) when \(H\) is nonempty. Under the conventional interpretation of finite simple undirected graphs, an \(F\)-free graph contains no subgraph isomorphic to \(F\); the copy need not be induced. These graph and containment conventions are not explicitly specified in the supplied record and require confirmation.

For graphs \(G\) and \(F\), set
\[
\mathcal H_F(G)=\{H:V(H)=V(G),\ E(H)\subseteq E(G),\ H\text{ is }F\text{-free}\}.
\]
Whenever this family is nonempty and \(G\) is finite and nonempty, define
\[
d_x(G,F)=\max_{H\in\mathcal H_F(G)}\delta(H).
\]
Thus \(d_x(G,F)\) is the largest minimum degree of an \(F\)-free spanning subgraph of \(G\). Let \(K_\Delta\) denote the complete graph on \(\Delta\) vertices. A graph is \(\Delta\)-regular if every vertex has degree \(\Delta\).

The supplied problem asks to characterize the graphs \(F\) for which
\[
\begin{gathered}
\exists\,\varepsilon>0\ \exists\,\Delta_0\in\Z_{\ge1}\ \forall\,\Delta\in\Z_{\ge1},\quad
\Delta\ge\Delta_0\ \Longrightarrow\\
\exists\text{ a finite nonempty }\Delta\text{-regular graph }G_\Delta
\quad\text{such that}\quad
d_x(G_\Delta,F)<(1-\varepsilon)d_x(K_\Delta,F).
\end{gathered}
\]
Here \(\varepsilon\) and \(\Delta_0\) may depend on \(F\), and \(G_\Delta\) may depend on \(F\) and \(\Delta\); no order or connectedness constraint on \(G_\Delta\) is supplied. The comparison uses \(K_\Delta\), whose degree is \(\Delta-1\), exactly as recorded.

The displayed interpretation is well-defined for finite \(F\) having at least one edge, since the edgeless spanning subgraph is then \(F\)-free. The source does not explicitly restrict \(F\) this way. If edgeless forbidden graphs are allowed, \(\mathcal H_F(G)\) can be empty, and no value for the corresponding maximum is supplied. Thus the full intended domain of \(F\), and the treatment of these boundary cases, remain unresolved; the conventional interpretation above is provisional.
\subsection{Short English statement}
Characterize the forbidden graphs \(F\) for which, for every sufficiently large integer \(\Delta\), some \(\Delta\)-regular graph has its best \(F\)-free spanning minimum degree smaller by a fixed positive proportion than the corresponding value for \(K_\Delta\). The graph conventions and edgeless forbidden-graph boundary cases require clarification.
\subsection{Sources}
\begin{itemize}
\item Supplied local UnsolvedMath registry record, permanent ID \(30003336\), source problem code \texttt{catalogue entry}, in the Oberwolfach Reports Open Problems collection. The supplied title is ``Regular Graphs Beating Clique Homomorphism Densities.'' No individual proposer is supplied; the proposed year is recorded as 2017.
\item Supplied primary-source citation: \emph{Combinatorics (2018)}, Oberwolfach Reports, DOI 10.4171/owr/2017/1, pp.\ 5--81. Original OWR report: \url{https://doi.org/10.4171/owr/2017/1}. The background describes these reports as documenting mathematical workshops and their open-problem discussions. The citation's 2018 label and the record's 2017 year information are retained without resolving their relationship.
\item Supplied literature reference: Foucaud--Krivelevich--Perarnau, ``large subgraphs without short cycles,'' \url{https://arxiv.org/abs/1401.4928}.
\item Literature assessment dated exactly \texttt{2026-08-21}: the supplied curation labels the problem open and states that a regular-graph homomorphism-density bound is proved in important cases, including selected even cycles, while an extracted conjecture for general \(F\) and a noted \(C_4\) case are not fully resolved. This is supplied dated curation derived from the OWR material, not an independently verified assessment. Its relation to the minimum-degree extremal statement above requires review.
\item The registry labels the statement \texttt{corrected\_verified} and reports that it was checked against the official Oberwolfach report and preserves its hypotheses, notation, and mathematical target. This verification claim is supplied provenance; no source was independently visited or checked for this draft.
\end{itemize}
\end{document}
